% You should *also* have a ND formatting guide to ensure that you have
% all the relevant parts, put the captions in the right place, etc.
% Just because you have this wonderful style classfile doesn't mean
% that it removes *all* the formatting onus from you.  :-)
% Although be warned that the Graduate School has been known to let
% their official formatting guide get out of date. When in doubt,
% the Microsoft Word example seemed to be the only thing kept
% consistently up-to-date in 2013, and is probably the safest thing
% to consult.d

% CTRL+B builds with latexmk 
% CTRL+SHIFT+B and then "Build With: Compile to PDF -Script Builder" emmulates the ./compile.sh script

\ProvidesFile{template.tex}
    [2016/10/16 v3a.2016%
     Template file for NDdiss2e class]
\documentclass[draft, numrefs, sort&compress, noinfo]{nddiss2e} %draft
             % One of the options draft, review, final must be chosen.
             % One of the options textrefs or numrefs should be chosen
             % to specify if you want numerical or ``author-date''
             % style citations.
             % Other available options are: 
             % 10pt/11pt/12pt (available with draft only)
             % twoadvisors
             % noinfo (sho uld be used when you compile the final time
             %         for formal submission)
             % sort (sorts multiple citations in the order that they're
             %       listed in the bibliography)
             % compress (compresses numerical citations, e.g. [1,2,3]
             %           becomes [1-3]; has no effect when used with
             %           the textrefs option)
             % sort&compress (sorts and compresses numerical citations;
             %           is identical to sort when used with textrefs)

\usepackage{placeins}   % for FloatBarrier
\usepackage[T1]{fontenc}
\usepackage{rotating}
% \usepackage{capt-of}  % using this broke my \ref{tab:triggers} table
% \usepackage{subfig}   % using this broke my \ref{tab:triggers} table
\usepackage{multirow}
\usepackage{threeparttable}
% \usepackage{booktabs}
\usepackage{makecell} 
\usepackage[caption=false]{subfig}
\captionsetup[subfigure]{margin=15pt} %{style=default, margin=0pt, parskip=0pt, hangindent=0pt, indention=0pt, singlelinecheck=true}
% \usepackage{subcaption}
\usepackage{lineno}
\usepackage{amssymb}
\usepackage{pifont} % http://ctan.org/pkg/pifont
\usepackage{lipsum}

\hypersetup{draft}

\begin{document}
\frontmatter                        % All the items before Chapter 1 go  in ``frontmatter''

% Line-breaks in the title have to be protected with `\protect`.
% \title{Gnus and you \protect\\ a Brief}
\title{}           % Title

\author{Christopher McGrady}              % Author's name
\work{Dissertation}                 % ``Dissertation'' or ``Thesis''
\degaward{Doctor of Philosophy}     % Degree you're aiming for.
                                    % Should be one of the following options:
                                    % ``Doctor of Philosophy'' (do NOT include ``in Subject'')
                                    % ``Master of Science \\ in \\ Subject''
\advisor{Kevin Lannon}              % Advisor's name
\program{Physics}                   % Name of the graduate program

\maketitle                          % The title page is created now

 % You must use either the \makecopyright option or the \makepublicdomain option.
 % \copyrightholder{}               % If you're not the copyright holder
 % \copyrightyear{}                 % If the copyright is not for the current year
\makecopyright                      % If not making your work public domain
                                    % uncomment out \makecopyright
 % \makepublicdomain                % Uncomment this to make your work public domain

 % Including an abstract is optional for a master's thesis, and required for a
 % doctoral dissertation.
% \include{abstract}                  % put it in a file to be included

 % Including a dedication is optional.
% \renewcommand{\dedicationname}{\mbox{}} % Replace \mbox{} if you want something else.
% \include{dedication}

\tableofcontents 
\listoffigures
\listoftables

 % Including a list of symbols is optional.
 %% \renewcommand{\symbolsname}{newsymname} % Replace ``newsymname'' with
                                            % the name you want, and uncomment
 % \begin{symbols}
 % \end{symbols}
 %                       % Use one of the two choices to add symbols text
 % \include{symbols}

 % Including a preface is optional.
 %% \renewcommand{\prefacename}{ } % If you want another Preface name, add
                                   % something else, and uncomment.
 % \begin{preface}
 % \end{preface}
 %                       % Use one of the two choices to add preface text
 % \include{preface}

 % Including an acknowledgements section may or may not be optional. It's hard to
 % tell from the information available in Spring 2013.
% \renewcommand{\acknowledgename}{ } % If you want another Acknowledgement name
                                       % add something else, and uncomment
 %                       % Use one of the two choices to add acknowledge text
% \include{acknowledgement}
\mainmatter
 % Place the text body here.
 % \include{chapter-one}
 % Begin each chapter with \chapter{Title}.

% Compile the document using pdflatex thesis.tex
\linenumbers

% If you have appendices, add them here.
% Begin each one with \chapter{TITLE} as before. The \appendix command takes
% care of renaming chapter headings and creates a new page in the Table of
% Contents for them.
% \appendix


\backmatter              % Place for bibliography and index

% \bibliographystyle{nddiss2e}
% \bibliography{bibliography}           % input the bib-database file name

\end{document}

%%
\endinput
%%
